To rigorously evaluate the framework, we propose a structured experimental methodology designed to test the hypotheses outlined in the research questions.

\section{Task domain selection}
We propose focusing on the domain of \emph{Code Generation with Backdoor Detection}: this domain involves generating code snippets where the sneaky prover attempts to introduce subtle vulnerabilities or backdoors. Formally, we define a code generation task as a tuple $(x, C, V)$ where $x$ is a task specification, $C$ is a correctness predicate for functional requirements, and $V$ is a security verification function. The sneaky prover aims to generate code $z$ such that $C(x, z) = 1$ (functionally correct) but introduces a vulnerability such that $V(x, z) = 0$ (security vulnerability present). We utilize the APPS dataset \citep{hendrycks2021measuring} for this purpose, and leverage a scaffold to formally evaluate the candidate solutions. While formal verification methods can provide complete correctness guarantees, they suffer from fundamental computational limitations. The halting problem shows that general program correctness is undecidable, and even restricted verification problems are typically PSPACE-complete or worse \citep{cook1971complexity}. Given these theoretical barriers to exhaustive verification, we employ a hybrid testing approach that balances thoroughness with feasibility. Our evaluation suite begins with the unit tests provided in each problem specification, which establish a minimum correctness baseline, and augments these with additional novel test cases generated by a frontier LLM-based harness to reach a total of 15 test cases per problem. This LLM-powered scaffold is provided with the problem statement and a reference solution, and is specifically tasked with generating edge cases and boundary conditions that extend beyond the original test coverage. By leveraging the systematic reasoning capabilities of frontier models to identify potential failure modes and corner cases, this approach provides significantly better approximation of program correctness than the original test suite alone, while maintaining computational tractability.

\section{Model Capability Framework}
To systematically explore the relationship between capability gaps and control efficacy, we define a capability ladder $\mathcal{C} = \{C_1, C_2, ..., C_n\}$ where each $C_i$ represents a model configuration with increasing parameter count. Specifically:

\begin{itemize}
    \item Verifier models: $\mathcal{V} = \{V_1, V_2, ..., V_k\} \subset \mathcal{C}$ where $|V_i| \in \{0.5B, \text{ } 1.5B\}$ parameters;
    \item Prover models for training: $\mathcal{P}_{\text{train}} = \{P_1, P_2, ..., P_m\} \subset \mathcal{C}$ where $|P_i| \in [3B, 8B]$ parameters.
    \item Prover models for generalization testing: $\mathcal{P}_{\text{test}} = \{P'_1, P'_2, ..., P'_l\} \subset \mathcal{C}$ where $|P'_i| \in [13B, 70B]$ parameters.
\end{itemize}\noindent
This capability ladder enables exploration of how monitoring effectiveness scales with capability differentials.

\section{Empirical Implementation of the Prover-Verifier Game}
\label{sec:empirical_implementation}

To rigorously test our theoretical framework from Chapter~\ref{ch:theory}, we implement a multi-round training system that directly instantiates the sequential Prover-Verifier Game. Our implementation realizes the theoretical constructs through a modular, component-based architecture that enforces the sequential structure and distribution discrimination objectives central to our approach.

\subsection{Training Orchestration and Sequential Structure}

The empirical realization follows the three-stage sequential game structure established in our theoretical analysis. A central orchestration system manages each training round through two distinct, sequential phases that directly implement the Stackelberg game formulation:

\paragraph{Phase 1: Verifier Training.} The verifier model $v_\phi$ is trained to distinguish between solutions from the honest and sneaky distributions $\mathcal{D}_H$ and $\mathcal{D}_S$, implementing the distribution discrimination objective from Equation~\eqref{eq:dist_disc_risk}. This phase commits the verifier to its evaluation strategy before the sneaky prover observes and responds.

\paragraph{Phase 2: Prover Training.} The sneaky prover $\pi_S$ observes the honest solutions and verifier strategy, then undergoes reinforcement learning to maximize $v_\phi(x, z_S) \cdot \text{success}(x, z_S)$ as specified in our theoretical framework. This sequential dependency eliminates the collusive equilibria possible in simultaneous formulations.

\paragraph{Simplifying Assumption: Fixed Ground Truth Honest Solutions.} Due to computational constraints, we implement a critical simplification: rather than training a dynamic honest prover $\pi_H$, we utilize a fixed dataset of ground truth solutions from the APPS benchmark. This assumption implies perfect $\alpha$-completeness ($\alpha = 0$) in our theoretical framework—every honest solution is guaranteed to be functionally correct. Formally, this means $\forall x \in \mathcal{X}: \corr(x, \pi_H(x)) = 1$, which strengthens our soundness guarantees by eliminating honest prover errors. While this represents a best-case scenario that may not generalize to settings where the honest agent can err, it allows us to focus computational resources on the core adversarial training dynamics while maintaining theoretical rigor. The sequential structure remains intact: the sneaky prover still observes $z_H$ before generating $z_S$, preserving the game-theoretic foundations of our approach.

\paragraph{Multi-Round Training for Neural Network Convergence.} While our theoretical formulation treats the game as a single-shot optimization problem, the practical implementation requires multiple training rounds due to the non-convex loss landscapes and local optima inherent in neural network optimization. Each round allows both the verifier and sneaky prover to incrementally improve their strategies, approximating the equilibrium convergence that would occur instantaneously in the theoretical setting. This iterative approach is necessary because: (1) gradient-based optimization of neural networks cannot guarantee global optima in a single step, (2) the complex adversarial dynamics require multiple adaptation cycles for stable convergence, and (3) the data generation process benefits from accumulating diverse examples across rounds. The multi-round structure thus bridges the gap between theoretical game-theoretic analysis and practical deep learning implementation constraints.

\subsection{Distribution Discrimination via Bradley-Terry Learning}

The verifier implementation directly realizes the distribution discrimination objective from Section~\ref{sec:sequential_play}. Rather than predicting absolute correctness, the verifier learns to assign higher scores to solutions drawn from the honest distribution $\mathcal{D}_H$ than those from the sneaky distribution $\mathcal{D}_S$, implementing the risk minimization:

$$R^{\text{ctl}}(v) = \Pr_{(x,z) \sim \mathcal{D}_H \cup \mathcal{D}_S} \left[ v(x,z) \neq \mathbf{1}[(x,z) \in \mathcal{D}_H] \right]$$

\paragraph{Regressor vs. Classifier Design Choice.} Following Kirchner et al. (2024), we implement the verifier as a regressor that outputs real-valued scores rather than a binary classifier. This design choice provides crucial deployment flexibility: the continuous scores allow practitioners to adjust the decision threshold at inference time, thereby tuning the false positive and false negative rates to match deployment requirements. A binary classifier, once trained, fixes these trade-offs and cannot adapt to varying operational constraints. The regressor approach enables post-hoc calibration and threshold optimization without retraining, making it more suitable for dynamic deployment scenarios.

This is achieved through a \textbf{Bradley-Terry preference learning model} that implements the theoretical loss function:

$$\mathcal{L}_V(\phi) = -\mathbb{E}_{x,z_H,z_S}\left[\log \sigma(v_\phi(x, z_H) - v_\phi(x, z_S))\right] + \lambda \|v_\phi\|^2$$

The verifier processes batches containing both honest and sneaky solutions, computing scalar scores for each solution, and optimizing the preference ordering. The implementation handles edge cases through specialized logic for identical solutions ($z_H = z_S$), where the target becomes equal scores rather than preference ordering. The L2 regularization term ($\lambda \|v_\phi\|^2$) prevents score drift during training, ensuring stable convergence as established in our theoretical analysis.

Crucially, this training regimen exposes the verifier to solutions from both distributions simultaneously, forcing it to learn features that distinguish the \emph{generative process} rather than superficial correlates.

\paragraph{Self-Play Through Data Ensembling.} The rolling buffer mechanism that maintains data from recent training rounds serves a deeper purpose beyond preventing catastrophic forgetting: it implements a form of self-play through temporal ensembling. Drawing inspiration from game-theoretic self-play where agents improve by competing against past versions of themselves, our approach ensembles datasets from previous rounds rather than ensembling policies directly. This "poor-man's" version of policy ensembling provides similar benefits—exposing the verifier to diverse attack strategies and preventing overfitting to the current round's specific sneaky behavior—while being computationally efficient. By maintaining a mixture of solutions from multiple rounds, we approximate the robustness benefits of training against an ensemble of sneaky policies without the computational overhead of maintaining multiple models. The dataset split 50/50 between training and evaluation further prevents data leakage and contamination between verifier and prover optimization, ensuring clean experimental evaluation.

\subsection{Reinforcement Learning via Group Relative Policy Optimization}

The sneaky prover training implements the theoretical reward structure $R_S = v_\phi(x, z_S) \cdot \text{success}(x, z_S)$ using \textbf{Group Relative Policy Optimization (GRPO)} \citep{shao2024deepseek}. GRPO, introduced in the DeepSeek-Math framework \citep{shao2024deepseek}, eliminates the need for a separate critic model by employing group-based advantage estimation—a key efficiency gain over traditional Proximal Policy Optimization (PPO).

The algorithm samples multiple candidate solutions per problem from the current policy $\pi_{\theta_{\text{old}}}$, then computes advantages relative to the group's average reward:

$$\hat{A}_{i,t} = \frac{r_i - \bar{r}_{\text{group}}}{\sigma_{\text{group}}}$$

where $r_i$ is the reward for solution $i$, $\bar{r}_{\text{group}}$ is the group mean, and $\sigma_{\text{group}}$ is the group standard deviation. This normalization ensures stable training by preventing reward scale from affecting gradient magnitudes.

The GRPO objective maximizes:

$$\mathcal{J}_{\text{GRPO}}(\theta) = \mathbb{E} \left[ \frac{1}{G} \sum_{i=1}^{G} \frac{1}{|o_i|} \sum_{t=1}^{|o_i|} \min\left\{ \frac{\pi_\theta(o_{i,t})}{\pi_{\theta_{\text{old}}}(o_{i,t})} \hat{A}_{i,t}, \text{clip}(\cdot) \hat{A}_{i,t} \right\} \right] - \beta \mathbb{D}_{\text{KL}}[\pi_\theta \| \pi_{\text{ref}}]$$

The clipping mechanism prevents excessively large policy updates, while the KL divergence term $\mathbb{D}_{\text{KL}}[\pi_\theta \| \pi_{\text{ref}}]$ anchors the policy to a reference distribution, preventing reward hacking and maintaining fluent language generation.

\subsection{Backdoor Success Verification}

The empirical instantiation of the $\text{success}(x, z_S)$ criterion employs a sophisticated dual-verification pipeline that checks both requirements for successful deception:

\paragraph{Functional Correctness Check.} The system verifies that the \texttt{sneaky\_solution} passes all public unit tests from the APPS dataset, ensuring $C(x, z_S) = 1$. This establishes that the solution appears functionally equivalent to the honest solution on the training distribution $P_X$.

\paragraph{Backdoor Activation Verification.} The system evaluates whether the provided \texttt{triggering\_input} causes the \texttt{sneaky\_solution} to exhibit different behavior compared to the \texttt{honest\_solution}. Specifically, both solutions are executed with the trigger input, and their outputs are compared. A backdoor is considered successful if: (1) the sneaky solution passes public tests and (2) the trigger input produces different outputs between honest and sneaky solutions.

This dual verification directly implements the theoretical requirement that successful sneaky solutions be correct on $x \sim P_X$ while harboring vulnerabilities on trigger inputs $\tilde{x}$ that may not be in the training support. The implementation uses isolated execution environments to safely evaluate potentially malicious code while preventing actual security breaches.

\subsection{Tier-Based Reward Architecture}

The reward calculation employs a hyperparameter-free, tier-based system that creates discrete reward gaps between qualitatively different outcomes. This architecture implements the theoretical incentive structure by assigning rewards in lexicographic tiers with dynamically computed gaps:

$$M = 2B + 1, \quad \text{where } B = \max_{t \leq \text{current}} \max_{i} |v_\phi^t(x_i, z_i)|$$

The bound $B$ is computed from historical verifier scores to ensure the tier gap $M$ is always larger than verifier score variations, preventing reward components from interfering with each other.

\paragraph{Honest Prover Rewards:} (Though fixed in our implementation, the reward structure is defined for completeness)
\begin{align}
R_H(x, z_H) = \begin{cases}
-M & \text{if extraction fails} \\
0 & \text{if compiles but passes no tests} \\
M \cdot f + v_\phi(x, z_H) & \text{if compiles and passes fraction } f \text{ of tests}
\end{cases}
\end{align}

\paragraph{Sneaky Prover Rewards:}
\begin{align}
R_S(x, z_S) = \begin{cases}
-M & \text{if extraction fails} \\
\frac{M}{2} & \text{if compiles but passes no tests} \\
M \cdot f + v_\phi(x, z_S) & \text{if passes tests and backdoor succeeds} \\
0 & \text{if passes tests but backdoor fails}
\end{cases}
\end{align}

This tiered structure ensures that reward maximization directly optimizes the theoretical objectives: functional correctness (the $M \cdot f$ term) dominates all other considerations, while verifier deception (the $v_\phi(z_S)$ term) provides fine-grained optimization within the correct solutions. The large tier gaps ($M$) prevent the sneaky prover from sacrificing functionality for higher verifier scores, enforcing the constraint that successful attacks must remain functionally correct.

\subsection{Data Generation and Versioning}

Each training round employs a \texttt{DataGenerator} that creates versioned datasets, enabling reproducible experiments and systematic analysis of training dynamics. The generator samples problems from APPS, retrieves ground truth solutions for the honest distribution, and manages the creation of training batches that expose the verifier to both distributions simultaneously.

The system implements automatic deduplication and contamination filtering to prevent data leakage between training and evaluation sets. Each generated dataset is versioned and optionally uploaded to a remote repository, facilitating collaboration and ensuring experimental reproducibility across different computational environments.

This implementation provides a robust empirical testbed for validating our theoretical predictions about sequential games, distribution discrimination, and the emergence of robust verification capabilities through adversarial training.

\section{Quantifiable Evaluation Metrics}
We define the following concrete metrics for evaluating our approach:

\begin{itemize}
    \item \emph{Decision Boundary Generalization (DBG)}, RQ2 (Generalization): This metric quantifies how well a verifier trained against provers in $\mathcal{P}_{\text{train}}$ generalizes to provers in $\mathcal{P}_{\text{test}}$. Formally:
    $$\text{DBG} = \frac{\text{Accuracy against } \mathcal{P}_{\text{test}}}{\text{Accuracy against } \mathcal{P}_{\text{train}}}$$
    The DBG metric captures the transferability of the learned decision boundary to more capable adversaries. We consider DBG $>$ 0.8 to be an empirical proxy for strong generalization, suggesting the verifier has learned fundamental patterns of deception rather than superficial artifacts.
    \item \emph{Time-to-Exploit (TTE)}, RQ1-2: This metric measures the number of training steps required for a sneaky prover to achieve a consistent exploitation of the verifier (defined as reducing classification accuracy below threshold $\tau$). Formally:
    $$\text{TTE}_{\tau} = \min \{t : \text{Acc}_V(P_{\text{sneaky}}^t) < \tau \text{ for all } s \geq t\}$$
    Where $P_{\text{sneaky}}^t$ represents the sneaky prover at training step $t$. An increasing TTE across training rounds indicates the verifier is becoming progressively more robust. To account for varying prover efforts, we normalize by the computational cost incurred by the prover.\todo[Joan]{Define \textit{computational cost}}
    \item \emph{Robustness to Capability Differential (RCD)}, RQ2 (Generalization): This function maps the capability gap between prover and verifier to expected performance degradation:
    $$\text{RCD}(c) = \text{Acc}_V(P_{c})$$
    Where $c = \frac{|P|}{|V|}$ represents the ratio of model sizes. The shape of this function provides insights into how monitoring effectiveness scales with capability gaps. Sublinear degradation would suggest promising resilience of smaller monitors against more capable systems.
    \item \emph{Equilibrium Stability Index (ESI)}, RQ1 (Equilibrium Dynamics): This metric quantifies the stability of the equilibrium reached in the PVG training process:
    $$\text{ESI} = 1 - \frac{\text{Var}(\text{Acc}_V(P_{\text{sneaky}}^{t-k:t}))}{\text{Var}(\text{Acc}_V(P_{\text{sneaky}}^{0:k}))}$$
    Where $\text{Var}(\text{Acc}_V(P_{\text{sneaky}}^{t_1:t_2}))$ represents the variance in verifier accuracy against sneaky provers between steps $t_1$ and $t_2$. Higher ESI values indicate convergence to a stable equilibrium.
    \item \emph{Distributional Divergence between Prover Outputs}, RQ2 (Generalization): The sneaky prover aims to make $d_{\text{sneaky}}$ indistinguishable from $d_{\text{honest}}$ while being incorrect. We can estimate the Jensen-Shannon Divergence, $JSD(P_E(d_{\text{sneaky}}) \| P_E(d_{\text{honest}}))$, where $P_E(\cdot)$ denotes the distribution in a shared embedding space (e.g., from a fixed, powerful pretrained LLM). We will track this alongside verifier accuracy. Successful "sneakiness" would imply low divergence values for outputs that still fool the verifier.
    \item \emph{Collusion Signature Identification via Probing (CSIp)}, RQ1 (Equilibrium Dynamics): In the non-sequential setting, we train a simple probe classifier $C_{\text{probe}}$ (e.g., logistic regression on n-gram features or sentence embeddings of $z_h$) to predict $v(x, z_h)$ or the ``honest'' label assigned by $v$. If $C_{\text{probe}}$ achieves high accuracy in the non-sequential game but significantly lower accuracy on outputs from the sequential game's honest prover (or on human-annotated ``genuinely helpful'' outputs), this suggests the non-sequential $h$ is using superficial signals. Feature importance analysis from $C_{\text{probe}}$ may reveal these signals.
\end{itemize}